\section{问题1：预热平衡阶段}

\subsection{模型建立}

预热平衡阶段采用附录2 的\textbf{常物性}，温度方程与水分方程分别为
\begin{equation}
  \rho c_p\frac{\partial T}{\partial t}
  =\frac{k}{r}\frac{\partial}{\partial r}\!\left(r\frac{\partial T}{\partial r}\right),
  \qquad
  \frac{\partial C}{\partial t}
  =\frac{1}{r}\frac{\partial}{\partial r}\!\left(D(C)\,r\frac{\partial C}{\partial r}\right),
\end{equation}
\begin{equation}
  D(C)=7\times10^{-9}\exp\!\left(-\frac{0.89}{C}\right)\quad[\mathrm{m^2/s}] .
\end{equation}

\textbf{结构性单向耦合}：$D$ \textbf{不含温度}，温度方程只通过``物性是否随 $C$ 变''
与水场耦合，而附录2 给的是常物性。故问题1 可严格拆成
\textbf{水分场独立求解 $\to$ 温度场代入已知 $C(r,t)$ 求解}。
这不是``为了省事''，而是方程结构本身允许的降维。

\textbf{初始条件与边界条件}：
\begin{equation}
  T(r,0)=28\ ^\circ\mathrm{C},\quad C(r,0)=2.55\ \mathrm{kg/kg};
\end{equation}
\begin{equation}
  r=0:\ \frac{\partial T}{\partial r}=\frac{\partial C}{\partial r}=0;
\end{equation}
\begin{equation}
  r=R_0:\ -k\frac{\partial T}{\partial r}\bigg|_{R}=h\left(T_s-T_\infty\right),
  \qquad
  -D\frac{\partial C}{\partial r}\bigg|_{R}=h_m\left(C_s-\Cinf\right).
\end{equation}

\subsection{数值方法}

\textbf{空间离散}：采用有限体积法（FVM），中心用\textbf{半控制体}
（半径 $\Delta r/2$ 的实心圆柱），天然规避 $1/r$ 奇点；
界面通量用\textbf{调和平均}，保证变 $D$ 时通量正确。

\textbf{网格}：均匀网格在表面早期欠分辨（$t=1$ s 时边界层厚度仅 $0.07$ mm，
均匀 $N=640$ 误差仍达 $4\times10^{-4}$）。故采用向表面加密的\textbf{渐变网格}
$N=200$、$\gamma=1.5$（$\Delta r=0.00707\sim0.14981$ mm），
以约一半计算量取得优于均匀 $N=640$ 的精度。

\textbf{时间推进}：半控制体 FVM 中心处的显式稳定性上限约 $\Delta r^{2}/(4\alpha)$，
均匀网格 $\Delta r=0.5$ mm 时为 $0.37$ s，无法满足 1 s 输出要求。
故采用 L-稳定的 \textbf{Backward Euler} 隐式推进，步长由误差控制自适应调整，
输出时刻精确落点。

\textbf{非线性迭代}：水分方程因 $D(C)$ 非线性，采用 \textbf{Picard 迭代}，
保持 M-矩阵性质，$C$ 的正性由构造保证（平均 2.92 次、最多 3 次）。
收敛采用尺度归一化的\textbf{双重判据}（更新量与真残差都要满足）。

\textbf{输出采样}：$r=0$ 用偶函数二次外推，$0<r<R_0$ 用 PCHIP 保形插值，
$r=R_0$ 用 Robin 边界重构值。

\begin{figure}[H]
  \centering
  \includegraphics[width=0.75\textwidth]{FIG-08_fvm.png}
  \caption{一维圆柱有限体积离散。(a) 横截面的环形控制体与中心半控制体；
           (b) 径向剖面，界面通量与表面 Robin 重构}
  \label{fig:FIG08}
\end{figure}

\subsection{结果}

\begin{figure}[H]
  \centering
  \includegraphics[width=0.78\textwidth]{G4_p1_evolution.png}
  \caption{问题1 温度场与含水率场的径向演化（0--1800 s）}
  \label{fig:G4}
\end{figure}

\textbf{温度场}：表面温度从 $28\ ^\circ\mathrm{C}$ 升至约 $36.8\ ^\circ\mathrm{C}$，
中心升至约 $33.6\ ^\circ\mathrm{C}$，表面始终高于中心，符合外部热风加热的物理过程。

\textbf{含水率场}：中心含水率在 1800 s 内基本不变（$2.55$ kg/kg），
表面从 $2.55$ 降至约 $1.51$ kg/kg，呈明显的\textbf{表面干燥}特征。

\input{tables/tab1}
\begin{table}[H]
  \centering
  \caption{30 分钟内药材的水分浓度（kg/kg）}
  \label{tab:p1C}
  \footnotesize
  % 数据源：同上「水分浓度」sheet
  \begin{tabular}{@{}lrrrrr@{}}
    \toprule
    时间/s & 0 cm & 0.5 cm & 1 cm & 1.5 cm & 2 cm \\
    \midrule
    100 & 2.5500 & 2.5500 & 2.5500 & 2.5500 & 2.2470 \\
    300 & 2.5500 & 2.5500 & 2.5500 & 2.5492 & 2.0508 \\
    600 & 2.5500 & 2.5500 & 2.5500 & 2.5353 & 1.8770 \\
    900 & 2.5500 & 2.5500 & 2.5497 & 2.5045 & 1.7547 \\
    1200 & 2.5500 & 2.5500 & 2.5482 & 2.4646 & 1.6586 \\
    1500 & 2.5500 & 2.5499 & 2.5445 & 2.4206 & 1.5787 \\
    1800 & 2.5500 & 2.5497 & 2.5383 & 2.3755 & 1.5103 \\
    \bottomrule
  \end{tabular}
\end{table}


完整结果见 \texttt{result1.xlsx}。

\subsection{分析与讨论}

\subsubsection*{（1）温度：表面升温快，中心严重滞后}

$t=1800$ s 时中心 $33.5763\ ^\circ\mathrm{C}$、表面 $36.7861\ ^\circ\mathrm{C}$，
\textbf{径向温差 $3.2098\ ^\circ\mathrm{C}$}，而中心仅完成烘房升温幅度的
$(33.58-28)/(41.51-28)=41.3\%$。

\begin{quote}
  这直接印证 $Bi_T=1.389$：\textbf{内外导热阻力相当，集总参数模型不可用。}
\end{quote}

\subsubsection*{（2）水分：干燥完全被限制在表层，中心未被触及}

\begin{table}[H]
  \centering
  \footnotesize
  \begin{tabular}{@{}lccc@{}}
    \toprule
    位置 & $t=1800$ s 含水率 & 变化 & 降幅 \\
    \midrule
    中心 & 2.5500 & $\mathbf{0.0000}$（30 min 内完全未变） & 0 \\
    表面 & 1.5103 & $-1.0073$ & $\mathbf{40.0\%}$ \\
    \bottomrule
  \end{tabular}
\end{table}

$C=2.5$ kg/kg 的等值线（``干燥前沿''）位于 $r\approx1.3$ cm——
即 30 min 内水分变化只发生在\textbf{外侧约 0.7 cm} 的薄层。

\subsubsection*{（3）干燥前沿呈``由表及里推进''形态}

从表 \ref{tab:p1C} 可见：表面持续单调下降；
而 $1.0$ cm 处直到约 $t=900$ s 才开始可见变化，
$1.5$ cm 处到 $t=1800$ s 才降 $0.17$。
这正是``\textbf{表面是唯一的门}、内部水分必须接力外迁''的直接体现。

\subsubsection*{（4）温度与水分的时间尺度差约 60 倍}

传热特征时间 $R_0^{2}/\alpha\approx2369$ s，
而传质特征时间 $R_0^{2}/D$（按 $D\approx5\times10^{-9}$ m$^2$/s）
$\approx8\times10^{4}$ s。实测表现完全一致：
\textbf{30 分钟内温度场已显著演化，而水分场只在表层动了 40\%}。

\begin{quote}
  \textbf{这解释了为什么烘干需要 2--3 天，而预热只需约 40 分钟。}
\end{quote}

\subsubsection*{（5）表面温度远未跟上烘房温度}

$t=1800$ s 时烘房已达 $41.513\ ^\circ\mathrm{C}$，而药材表面仅 $36.7861\ ^\circ\mathrm{C}$
——即表面只完成烘房升温幅度的 $65.0\%$。

\begin{quote}
  \textbf{注意}：这是 M0（\textbf{不含蒸发潜热项}）的结果。
  若计入蒸发吸热，表面温度还会更低。
  \textbf{M1 扩展不得直接套用本结论}，必须先在通量闭合的前提下重算。
\end{quote}
